\documentclass[12pt]{article}

\usepackage[utf8]{inputenc}
\usepackage[T1]{fontenc}
\usepackage{amsmath, amssymb}
\usepackage{graphicx}
\usepackage{float}
\usepackage{listings}
\usepackage{xcolor}
\usepackage[top=1in, bottom=1in, left=1in, right=1in]{geometry}
\usepackage{tikz}
\usetikzlibrary{matrix, arrows.meta, calc}
\usepackage{titling}
\setlength{\droptitle}{-3em}

\lstset{
    language=Python,
    basicstyle=\ttfamily\small,
    keywordstyle=\color{blue},
    commentstyle=\color{green!60!black},
    stringstyle=\color{red},
    numbers=left,
    numberstyle=\tiny\color{gray},
    frame=single,
    backgroundcolor=\color{gray!10},
    breaklines=true,
    showstringspaces=false
}

\lstdefinestyle{plainschema}{
    language=,
    basicstyle=\ttfamily\small,
    frame=single,
    backgroundcolor=\color{gray!10},
    breaklines=true,
    showstringspaces=false,
    numbers=none
}

\lstdefinestyle{sqlstyle}{
    language=SQL,
    basicstyle=\ttfamily\small,
    keywordstyle=\color{blue},
    commentstyle=\color{green!60!black},
    stringstyle=\color{red},
    frame=single,
    backgroundcolor=\color{gray!10},
    breaklines=true,
    showstringspaces=false,
    numbers=none
}

\geometry{margin=1in}

\newcommand{\schema}[1]{%
  \begin{center}%
  \setlength{\fboxsep}{8pt}%
  \colorbox{gray!10}{%
    \begin{minipage}{0.93\linewidth}\ttfamily\small\raggedright #1\end{minipage}}%
  \end{center}%
}

\title{HW1}
\author{
Greyson Knippel \\[4pt]
CPSC 321 --- Gonzaga University \\[4pt]
compiled with \LaTeX{} and Ti\emph{k}Z :)
}
\date{\today}

\begin{document}

\maketitle

\section*{part A}

\subsection*{question 1}

\schema{
AIRPORT(\underline{air\_id}, air\_name, city, state)\\[4pt]
FK: none
}

\texttt{air\_id} is the primary key since its unique, and minimal.

\begin{table}[H]
\centering
\begin{tabular}{|l|l|l|l|}
\hline
\textbf{air\_id} & \textbf{air\_name} & \textbf{city} & \textbf{state} \\
\hline
GEG & Spokane International & Spokane & WA \\
SEA & Seattle-Tacoma International & Seattle & WA \\
LAX & Los Angeles International & Los Angeles & CA \\
\hline
\end{tabular}
\end{table}

\subsection*{question 2}

\schema{
AIRLINE(\underline{air\_code}, airline\_name, air\_id, founded\_year)\\[4pt]
FK: air\_id -> AIRPORT(air\_id)
}

\texttt{air\_code} is the primary key since its unique, and minimal.
\texttt{air\_id} is a foreign key referencing the hub airport's \texttt{air\_id}
in AIRPORT.

\begin{table}[H]
\centering
\begin{tabular}{|l|l|l|l|}
\hline
\textbf{air\_code} & \textbf{airline\_name} & \textbf{air\_id} & \textbf{founded\_year} \\
\hline
AS & Alaska Airlines & SEA & 1932 \\
WN & Southwest Airlines & LAX & 1967 \\
UA & United Airlines & GEG & 1926 \\
\hline
\end{tabular}
\end{table}

\subsection*{question 3}

\schema{
FLIGHT\_ROUTE(\underline{air\_code}, \underline{flight\_num}, dep\_id, arr\_id,
distance)\\[4pt]
FK: air\_code -> AIRLINE(air\_code)\\
FK: dep\_id -> AIRPORT(air\_id)\\
FK: arr\_id -> AIRPORT(air\_id)
}

primary key is composite. neither attribute is unique on its own: an
airline operates many routes, so \texttt{air\_code} repeats, and two airlines
can both use flight number 122, so \texttt{flight\_num} repeats. together they
are unique, and removing either one breaks uniqueness, so the pair is minimal.

\texttt{dep\_id} and \texttt{arr\_id} are both foreign keys into AIRPORT.

\begin{table}[H]
\centering
\begin{tabular}{|l|l|l|l|l|}
\hline
\textbf{air\_code} & \textbf{flight\_num} & \textbf{dep\_id} & \textbf{arr\_id} & \textbf{distance} \\
\hline
AS & 122 & GEG & SEA & 224 \\
AS & 345 & SEA & LAX & 954 \\
WN & 122 & LAX & GEG & 1046 \\
\hline
\end{tabular}
\end{table}

\subsection*{question 4}

\schema{
PASSENGER(\underline{pass\_id}, first\_name, last\_name, email)\\[4pt]
Additional candidate key: \{email\}\\
FK: none
}

\texttt{pass\_id} and \texttt{email} are both candidate keys since both are
stated to be unique. \texttt{pass\_id} should be the primary key because
email addresses can change, and that would be a problem.

\begin{table}[H]
\centering
\begin{tabular}{|l|l|l|l|}
\hline
\textbf{pass\_id} & \textbf{first\_name} & \textbf{last\_name} & \textbf{email} \\
\hline
1 & Ana & Ruiz & aruiz@example.com \\
2 & Ben & Cho & bcho@example.com \\
3 & Dee & Patel & dpatel@example.com \\
\hline
\end{tabular}
\end{table}

\schema{
BOOKING(\underline{pass\_id}, \underline{air\_code}, \underline{flight\_num},
\underline{travel\_date}, seat, price)\\[4pt]
FK: pass\_id -> PASSENGER(pass\_id)\\
FK: (air\_code, flight\_num) -> FLIGHT\_ROUTE(air\_code, flight\_num)
}

the constraint ``booked once per route per date'' 
makes those four attributes the key. every removal of the attributes breaks it as well.

\texttt{(air\_code, flight\_num)} is a composite foreign key the parent key
is composite, so the reference has to carry both columns together.
\texttt{seat} and \texttt{price} are irrelevant to this table's keys.

\begin{table}[H]
\centering
\begin{tabular}{|l|l|l|l|l|l|}
\hline
\textbf{pass\_id} & \textbf{air\_code} & \textbf{flight\_num} & \textbf{travel\_date} & \textbf{seat} & \textbf{price} \\
\hline
1 & AS & 122 & 2026-10-02 & 14A & 189.50 \\
1 & AS & 122 & 2026-11-15 & 08C & 210.00 \\
2 & AS & 122 & 2026-10-02 & 14B & 205.75 \\
\hline
\end{tabular}
\end{table}

\section*{part B}

\subsection*{question 5}

\begin{center}
\begin{tikzpicture}[
  reltable/.style={
    matrix of nodes,
    ampersand replacement=\&,
    column sep=-\pgflinewidth,
    nodes={draw, rectangle, font=\ttfamily\scriptsize, minimum height=6.5mm,
           inner xsep=4pt, anchor=center}
  },
  relname/.style={font=\bfseries\footnotesize, anchor=east},
  fk/.style={-{Stealth[length=2mm]}, semithick, draw=blue!60!black, rounded corners=2pt}
]

\matrix (airport) [reltable, anchor=north west] at (0,0) {
  \underline{air\_id} \& air\_name \& city \& state \\
};
\matrix (airline) [reltable, anchor=north west] at (0,-2.2) {
  \underline{air\_code} \& airline\_name \& air\_id \& founded\_year \\
};
\matrix (fr) [reltable, anchor=north west] at (0,-4.4) {
  \underline{air\_code} \& \underline{flight\_num} \& dep\_id \& arr\_id \& distance \\
};
\matrix (bk) [reltable, anchor=north west] at (0,-6.6) {
  \underline{pass\_id} \& \underline{air\_code} \& \underline{flight\_num} \&
  \underline{travel\_date} \& seat \& price \\
};
\matrix (pax) [reltable, anchor=north west] at (0,-8.8) {
  \underline{pass\_id} \& first\_name \& last\_name \& email \\
};

\node[relname] at ([xshift=-3mm]airport.west) {AIRPORT};
\node[relname] at ([xshift=-3mm]airline.west) {AIRLINE};
\node[relname] at ([xshift=-3mm]fr.west) {FLIGHT\_ROUTE};
\node[relname] at ([xshift=-3mm]bk.west) {BOOKING};
\node[relname] at ([xshift=-3mm]pax.west) {PASSENGER};

\coordinate (chA) at ([xshift=7mm]bk.east);
\coordinate (chB) at ([xshift=15mm]bk.east);
\coordinate (topA) at ([yshift=7mm]airport.north east);
\coordinate (topB) at ([yshift=14mm]airport.north east);
\coordinate (dA) at ([yshift=4mm]fr-1-3.north);
\coordinate (dB) at ([yshift=9mm]fr-1-4.north);
\coordinate (e1) at ([xshift=-4mm]airport-1-1.north);
\coordinate (e2) at ([xshift=4mm]airport-1-1.north);

\draw[fk] (airline-1-3.north) -- ++(0,6mm) -| (airport-1-1.south);

\draw[fk] (fr-1-1.north) -- ++(0,7mm) -| (airline-1-1.south);

\draw[fk] (fr-1-3.north) -- (dA) -- (chA |- dA) -- (chA |- topA)
          -- (e1 |- topA) -- (e1);

\draw[fk] (fr-1-4.north) -- (dB) -- (chB |- dB) -- (chB |- topB)
          -- (e2 |- topB) -- (e2);

\draw[fk] (bk-1-2.north) -- ++(0,5mm) -| (fr-1-1.south);
\draw[fk] (bk-1-3.north) -- ++(0,10mm) -| (fr-1-2.south);

\draw[fk] (bk-1-1.south) -- ++(0,-6mm) -| (pax-1-1.north);

\end{tikzpicture}
\end{center}

the underlined attributes are primary keys. each arrow runs from a foreign key to
the primary key it references. the two arrows from BOOKING into FLIGHT\_ROUTE
are a single composite foreign key: \texttt{air\_code} and
\texttt{flight\_num} must be carried together, since the key they reference is
composite. AIRPORT is referenced three times: once by AIRLINE for its hub, and
twice by FLIGHT\_ROUTE for the departure and arrival airports.

\newpage

\section*{part C}

\subsection*{question 6}

\subsubsection*{(a)}

\schema{
MaintenanceRecord(\underline{aircraft\_id}, \underline{maintenance\_date},
\underline{mechanic\_id},\\
\hspace*{2em}aircraft\_model, manufacturer, mechanic\_name, work\_performed,\\
\hspace*{2em}parts\_cost, mechanic\_cert\_level)\\[4pt]
FK: none
}

the key means one row is ``this mechanic worked on this aircraft on this
date,'' so every other attribute should be a fact about that event.

the redundant
values are \textbf{\texttt{mechanic\_name}} and
\textbf{\texttt{mechanic\_cert\_level}}. both are facts about the mechanic, not
about a maintenance job, so they repeat on every record that mechanic performs:

\begin{table}[H]
\centering
\begin{tabular}{|l|l|l|}
\hline
\textbf{attribute} & \textbf{really a fact about} & \textbf{so it repeats} \\
\hline
\texttt{aircraft\_model} & the aircraft & every record for that aircraft \\
\texttt{manufacturer} & the aircraft (and its model) & every record for that aircraft \\
\texttt{mechanic\_name} & the mechanic & every record by that mechanic \\
\texttt{mechanic\_cert\_level} & the mechanic & every record by that mechanic \\
\texttt{work\_performed} & the visit itself & never \\
\texttt{parts\_cost} & the visit itself & never \\
\hline
\end{tabular}
\end{table}

a 737-900 is always built by Boeing, so the same value also repeats across
completely different aircraft that happen to share a model.

only \texttt{work\_performed} and \texttt{parts\_cost} change from one
maintenance visit to the next, so they are the only attributes that actually
belong in this table.

\begin{table}[H]
\centering
\footnotesize
\begin{tabular}{|l|l|l|l|l|l|l|l|l|}
\hline
\textbf{aircraft} & \textbf{date} & \textbf{mech} & \textbf{model} &
\textbf{manuf.} & \textbf{mech\_name} & \textbf{work} & \textbf{cost} &
\textbf{cert} \\
\hline
N101AS & 2026-03-04 & M7 & 737-900 & Boeing & Dana Weiss & brake pads & 1450.00 & A\&P \\
N101AS & 2026-06-18 & M7 & 737-900 & Boeing & Dana Weiss & hydraulics & 0.00 & A\&P \\
N220WN & 2026-06-18 & M7 & 737-700 & Boeing & Dana Weiss & oil change & 310.00 & A\&P \\
N101AS & 2026-09-02 & M3 & 737-900 & Boeing & Omar Diaz & tires & 2100.00 & IA \\
\hline
\end{tabular}
\end{table}

``Dana Weiss'' and ``A\&P'' are written three times because M7 appears three
times, and ``Boeing'' is written on all four rows despite two different
aircraft and two different models.

\newpage

\subsubsection*{(b)}

each group of facts from part (a) moves into its own table, and the keys are whatever
decides it. what is left behind in the original table is just the id, which
links back to the new table:

\schema{
AIRCRAFT\_MODEL(\underline{aircraft\_model}, manufacturer)\\
FK: none\\[8pt]
AIRCRAFT(\underline{aircraft\_id}, aircraft\_model)\\
FK: aircraft\_model -> AIRCRAFT\_MODEL(aircraft\_model)\\[8pt]
MECHANIC(\underline{mechanic\_id}, mechanic\_name, mechanic\_cert\_level)\\
FK: none\\[8pt]
MAINTENANCE\_RECORD(\underline{aircraft\_id}, \underline{maintenance\_date},
\underline{mechanic\_id},\\
\hspace*{2em}work\_performed, parts\_cost)\\
FK: aircraft\_id -> AIRCRAFT(aircraft\_id)\\
FK: mechanic\_id -> MECHANIC(mechanic\_id)
}

the key of MAINTENANCE\_RECORD does not change --- it is still the same three
attributes --- but \texttt{aircraft\_id} and \texttt{mechanic\_id} are now
foreign keys as well as key attributes. only \texttt{work\_performed} and
\texttt{parts\_cost} stay behind, since those are the only values that actually
describe the visit.

the fourth table is what removes the last of the redundancy. if
\texttt{manufacturer} were left in AIRCRAFT, ``Boeing'' would still be written
once per aircraft. giving the model its own table means each manufacturer is
recorded exactly once.

\begin{table}[H]
\centering
\footnotesize
\begin{tabular}{|l|l|}
\hline
\textbf{aircraft\_model} & \textbf{manufacturer} \\
\hline
737-900 & Boeing \\
737-700 & Boeing \\
A320-200 & Airbus \\
\hline
\end{tabular}
\qquad
\begin{tabular}{|l|l|}
\hline
\textbf{aircraft\_id} & \textbf{aircraft\_model} \\
\hline
N101AS & 737-900 \\
N220WN & 737-700 \\
N455UA & A320-200 \\
\hline
\end{tabular}
\end{table}

\begin{table}[H]
\centering
\footnotesize
\begin{tabular}{|l|l|l|}
\hline
\textbf{mechanic\_id} & \textbf{mechanic\_name} & \textbf{mechanic\_cert\_level} \\
\hline
M7 & Dana Weiss & A\&P \\
M3 & Omar Diaz & IA \\
M9 & Lena Petrov & A\&P \\
\hline
\end{tabular}
\end{table}

\begin{table}[H]
\centering
\footnotesize
\begin{tabular}{|l|l|l|l|l|}
\hline
\textbf{aircraft\_id} & \textbf{maintenance\_date} & \textbf{mechanic\_id} &
\textbf{work\_performed} & \textbf{parts\_cost} \\
\hline
N101AS & 2026-03-04 & M7 & brake pads & 1450.00 \\
N101AS & 2026-06-18 & M7 & hydraulics & 0.00 \\
N220WN & 2026-06-18 & M7 & oil change & 310.00 \\
N101AS & 2026-09-02 & M3 & tires & 2100.00 \\
\hline
\end{tabular}
\end{table}

these have the same information as the original table but are just more optimized.

\end{document}